\section{Homomorphic Encryption}

For an algorithm to be \emph{oblivious} or \emph{data-independent}, it must perform the same operations regardless of the data being processed. In other words, \gls{fhe} is \emph{branchless programming}; programming without conditional statements such as \texttt{if}, \texttt{switch}, \texttt{match}, and so on. (Conditional with respect to the data, the program can still be stateful etc.) Therefore, it is often natural to think of \gls{fhe} algorithms as logical/boolean circuits. 

\begin{itemize}
    \item \textbf{Homomorphic Encryption} allows computation on encrypted data
    \item \textbf{\gls{fhe}} can evaluate \emph{any} circuit consisting of addition $(+)$ and multiplication $(\cdot)$
    \item \textbf{Levelled \gls{fhe}} can only evaluate circuits of a certain depth, but are faster than their non-levelled counterparts
    \item \textbf{Somewhat HE} supports a limited number of operations before noise renders data unreadable, whereas FHE employs \emph{bootstrapping} to refresh noisy data, enabling unlimited, arbitrary computations
\end{itemize}

\noindent\textbf{Relevant resources:}
\begin{itemize}
    \item \cite{cryptoeprint:2025/860}
    \item \cite{cryptoeprint:2022/915}
    \item \cite{cryptoeprint:2021/435}
    \item \cite{cryptoeprint:2016/870}
    \item \cite{cryptoeprint:2016/421}
    \item \cite{cryptoeprint:2013/340}
    \item \cite{cryptoeprint:2012/144}
    \item \cite{cryptoeprint:2012/078}
    \item \cite{cryptoeprint:2011/277}
    \item \cite{arxiv:2022/textbook}
    \item \cite{park:sortinghat:2022}
    \item \cite{chillotti:onionring:2019}
\end{itemize}

\newpage

\subsection{Basics}

A tuple of setup, keygen, eval and decrypt...

\subsubsection{Encoding}



\subsubsection{Embedding}
DCRT!


\subsection{Hardness Assumptions}

Reductions from hard lattice problems: IVP, SVP, Gap versions, Shortest Integer Solution, LWE.

\subsection{Ciphertexts}

\gls{rlwe} ciphertext:
\begin{equation}
    \mathbf{RLWE}_{s}(m):\quad (a, b), \quad b = a\cdot s + \delta\cdot m + \varepsilon
\end{equation}

\gls{rgsw} ciphertext, defined in \cite{cryptoeprint:2013/340}:
\begin{equation}
    \mathbf{RGSW}_s(m):\quad ???
\end{equation}

An important operation for this thesis is defined between these two forms through \eqref{eq:externalprod}, where the inverse gadget matrix $G^{-1}(\cdot)$ satisfies \eqref{eq:externalprod_gadget}. This operation is called the \emph{external product} and is denoted by the $\boxdot$ operator.
\begin{align}
    \label{eq:externalprod}
    a \boxdot B &= G^{-1}(a)\cdot B \\
    \label{eq:externalprod_gadget}
    G^{-1}(a)\cdot G &= a + e \pmod{q}
\end{align}

A key observation is that an \gls{rgsw} ciphertext encodes a collection of \gls{rlwe} samples. It follows that we can convert, or expand, an \gls{rlwe} ciphertext to an \gls{rgsw} representation.

\begin{algorithm}
\caption{HomExpand}
\label{alg:HomExpand}
\begin{algorithmic}[1] % [1] adds line numbering every line
    \State \textbf{Input:} A set $\mathcal{S}$ of \gls{rlwe} ciphertexts
    \State \textbf{Output:} An \gls{rgsw} ciphertext
    % \If{HSIC$_{X \to Y} <$ HSIC$_{Y \to X}$}
    %     \State The causal direction is $X \to Y$
    % \ElsIf{HSIC$_{X \to Y} >$ HSIC$_{Y \to X}$}
    %     \State The causal direction is $Y \to X$
    % \Else
    %     \State Direction is ambiguous
    % \EndIf
    \For{$c = (a,b)$ in $\mathcal{S}$}
        \State Process data
    \EndFor
    \State Return result
\end{algorithmic}
\end{algorithm}

\clearpage

% \subsection{Ciphertexts (AI preview)}

% \gls{rlwe} ciphertext:
% \begin{equation}
%     \mathbf{RLWE}_{s}(m):\quad (a, b), \quad b = a\cdot s + \delta\cdot m + \varepsilon
% \end{equation}

% \gls{rgsw} ciphertext, defined in \cite{cryptoeprint:2013/340}:
% \begin{equation}\label{eq:rgsw}
%     \mathbf{RGSW}_s(m):\quad
%     \mathbf{C} \in \mathcal{R}_q^{2\ell \times 2},\qquad
%     \mathbf{C} = \mathbf{Z} + m\cdot\mathbf{G}
% \end{equation}
% Here $\mathbf{Z}\in\mathcal{R}_q^{2\ell\times 2}$ is a matrix whose rows are fresh
% $\mathbf{RLWE}_s(0)$ samples, $\mathbf{G} = \mathbf{I}_2\otimes\mathbf{g}
% \in\mathcal{R}_q^{2\ell\times 2}$ is the gadget matrix, and
% $\mathbf{g}=(1,\beta,\dots,\beta^{\ell-1})^{\!\top}$ for decomposition base $\beta$
% and level count $\ell$.

% An important operation is defined between these two forms through
% \eqref{eq:externalprod}, where the inverse gadget matrix $G^{-1}(\cdot)$ satisfies
% \eqref{eq:externalprod_gadget}.  This operation is called the \emph{external product}
% and is denoted by the $\boxdot$ operator.
% \begin{align}
%     \label{eq:externalprod}
%     a \boxdot B &= G^{-1}(a)\cdot B \\
%     \label{eq:externalprod_gadget}
%     G^{-1}(a)\cdot G &= a + e \pmod{q}
% \end{align}
% Applying the external product to an $\mathbf{RLWE}_s$ ciphertext
% $c = (a,b)$ and an $\mathbf{RGSW}_s(m)$ ciphertext $\mathbf{C}$ gives
% \begin{equation}\label{eq:ext_prod_result}
%     c \boxdot \mathbf{C}
%     = G^{-1}(a,b)\cdot\bigl(\mathbf{Z}+m\cdot\mathbf{G}\bigr)
%     \approx m\cdot(a,b),
% \end{equation}
% where the approximation follows because $G^{-1}(a,b)\cdot\mathbf{Z}$ contributes
% only a small error (each row of $\mathbf{Z}$ decrypts to zero) and
% $G^{-1}(a,b)\cdot\mathbf{G}\approx(a,b)$ by \eqref{eq:externalprod_gadget}.
% The result $m\cdot(a,b)$ is itself an $\mathbf{RLWE}_s$ ciphertext, enabling
% homomorphic scalar multiplication in $\mathcal{R}_q$.

% A key observation is that an \gls{rgsw} ciphertext encodes a collection of
% \gls{rlwe} samples---its $2\ell$ rows each constitute a pair in $\mathcal{R}_q^2$.
% It follows that we can convert, or \emph{expand}, a structured collection of
% \gls{rlwe} ciphertexts into an \gls{rgsw} representation, as described in
% Algorithm~\ref{alg:HomExpand}.

% \begin{algorithm}
% \caption{HomExpand}
% \label{alg:HomExpand}
% \begin{algorithmic}[1]
%     \State \textbf{Input:}  A set
%            $\mathcal{S} = \bigl\{c_j=(a_j,b_j)\bigr\}_{j=1}^{2\ell}$
%            of \gls{rlwe} ciphertexts under an auxiliary key $s'$, where $c_j$
%            encrypts $\bigl\lfloor q/\beta^{(j-1)\bmod\ell}\bigr\rfloor\cdot m$
%     \State \textbf{Output:} An \gls{rgsw} ciphertext
%            $\mathbf{C}\in\mathcal{R}_q^{2\ell\times 2}$
%     \State Initialise $\mathbf{C}\leftarrow\mathbf{0}^{2\ell\times 2}$,
%            $\;j\leftarrow 1$
%     \For{$c=(a,b)\in\mathcal{S}$}
%         \State $(a,b)\leftarrow\mathsf{KS}_{s'\to s}(a,b)$
%                \Comment{Re-encrypt under target key $s$}
%         \State $(a,b)\leftarrow\mathsf{ModSwitch}_{\,q'\to q}(a,b)$
%                \Comment{Reduce noise via modulus switching}
%         \State $\mathbf{C}[j,\,:\,]\leftarrow(a,\,b)$
%                \Comment{Assign as $j$-th row of RGSW matrix}
%         \State $j\leftarrow j+1$
%     \EndFor
%     \State \Return $\mathbf{C}$
% \end{algorithmic}
% \end{algorithm}

\clearpage
